\clearpage
\fancyhead[L]{\small KOMET.}
\hypertarget{native-442-part-lower}{}
\begin{center}
{\large KOMET.}

\noindent\rule{\linewidth}{0.4pt}

Göttingische gelehrte Anzeigen. Nr. 78. S. 769..771. 17. Mai 1821.

\noindent\rule{\linewidth}{0.4pt}
\end{center}

Wir haben bisher damit gezögert, sämtliche Beobachtungen des diesjährigen Kometen mitzuteilen, die an unserer Sternwarte gemacht wurden. Wir hofften nämlich, dass sich die darauf beruhenden Ergebnisse durch viel spätere Beobachtungen von anderen Orten noch verbessern ließen. Da sich diese Hoffnung nicht erfüllt hat, teilen wir nun hier sämtliche hiesigen Beobachtungen mit.

\Needspace{14\baselineskip}
\SourceLeaf{443}
\begin{center}
\begin{tabular}{lr|r|r|r}
\multicolumn{2}{c|}{1821} & Mittlere Zeit & Rektaszension & Nördl. Dekl. \\
\hline
Jan. & 30 & $7^{\mathrm{u}}\;34^{\mathrm{m}}\;32^{\mathrm{s}}$ & $359^\circ\;27'\;7''$ & $16^\circ\;4'\;36''$ \\
Febr. & 3 & $7\quad3\quad56$ & $359\quad3\quad54$ & $15\quad46\quad3$ \\
& 7 & $6\quad42\quad38$ & $358\quad45\quad5$ & $15\quad29\quad49$ \\
& 9 & $6\quad42\quad30$ & $358\quad36\quad24$ & $15\quad21\quad22$ \\
& 10 & $6\quad52\quad27$ & $358\quad32\quad19$ & $16\quad17\quad50$ \\
& 11 & $7\quad12\quad0$ & $358\quad28\quad27$ & $15\quad14\quad26$ \\
März & 1 & $7\quad18\quad7$ & $357\quad13\quad34$ & $14\quad8\quad37$ \\
& 5 & $7\quad5\quad28$ & $356\quad54\quad11$ & $13\quad43\quad5$ \\
\end{tabular}
\end{center}

\Needspace{14\baselineskip}
Herr von \textsc{Staudt}, von dessen hervorragendem Geschick im astronomischen Rechnen wir schon öfter Beispiele mitgeteilt haben, stützt auf diese Beobachtungen die folgenden \emph{parabolischen Elemente}.

\begin{center}
\small
\begin{tabular}{p{.73\linewidth}r}
Zeit des Durchgangs durchs Perihel 1821 März \mbox{.\,.\,.\,.\,.\,.\,.\,.\,.} & $21.61890$ \\
Logarithmus des kleinsten Abstands von der Sonne Meridian von Göttingen & $8.9645990$ \\
Argument der Breite im Perihel \mbox{.\,.\,.\,.\,.\,.\,.\,.\,.\,.\,.\,.\,.} & $169^\circ\;10'\;9''3$ \\
Länge des aufsteigenden Knotens \mbox{.\,.\,.\,.\,.\,.\,.\,.\,.\,.\,.\,.\,.\,.\,.\,.} & $48\quad44\quad14.7$ \\
Neigung der Bahn \mbox{.\,.\,.\,.\,.\,.\,.\,.\,.\,.\,.\,.\,.\,.\,.\,.\,.\,.\,.\,.\,.} & $106\quad40\quad16.4$ \\
\end{tabular}
\end{center}

\Needspace{17\baselineskip}
Diese Elemente geben sämtliche uns bekannt gewordenen Beobachtungen so gut wieder, dass keine weitere Verbesserung und noch weniger eine Bestimmung der Elliptizität möglich ist, solange keine viel späteren guten Beobachtungen von anderen Orten bekannt werden. Aus Platzgründen können wir hier nur den Vergleich der Elemente mit den hiesigen Beobachtungen beifügen.

\begin{center}
\begin{tabular}{lr|r|r}
& & \multicolumn{2}{c}{Unterschied.} \\
& & Rektasz. & Deklination \\
\hline
Januar & 30 & $-18''8$ & $+34''9$ \\
Februar & 3 & $+3.9$ & $-6.0$ \\
& 7 & $-13.1$ & $-36.0$ \\
& 9 & $-0.1$ & $+14.7$ \\
& 10 & $+2.6$ & $+7.3$ \\
& 11 & $-1.2$ & $-3.5$ \\
März & 1 & $-0.3$ & $-0.3$ \\
& 5 & $-5.3$ & $-20.6$ \\
\end{tabular}
\end{center}

\Needspace{11\baselineskip}
Dieser Komet unterscheidet sich von anderen durch seine starke Annäherung an die Sonne und musste daher um die Zeit des Perihels eine sehr große Helligkeit erreichen. Unter der Annahme, dass das Licht des Kometen geborgtes Sonnenlicht ist, ergab die Rechnung von Herrn \textsc{von Staudt} die folgenden Helligkeitswerte.

\begin{center}
\begin{tabular}{lr|r}
Januar & 21 & $0.139$ \\
Februar & 10 & $0.258$ \\
März & 1 & $0.832$ \\
& 5 & $1.253$ \\
\end{tabular}
\end{center}
\Needspace{7\baselineskip}
\begin{center}
\begin{tabular}{lr|r}
März & 18 & $19.588$ \\
& 19 & $32.519$ \\
& 20 & $59.932$ \\
& 21 & $110.610$ \\
\end{tabular}
\end{center}
\Needspace{7\baselineskip}
\begin{center}
\begin{tabular}{lr|r}
März & 22 & $128.139$ \\
& 23 & $81.238$ \\
& 24 & $46.005$ \\
& 25 & $27.980$ \\
\end{tabular}
\end{center}

Die Zahlen ab dem 18. März gelten für die Kulminationszeit in Göttingen. In den ersten Märztagen war der Komet sehr leicht mit bloßem Auge zu sehen und war fast so hell wie ein Stern dritter Größe. Deshalb bestand große Hoffnung, dass sich der Komet zur Zeit seiner größten Helligkeit trotz seiner Stellung nahe der Sonne tagsüber mit lichtstarken Instrumenten im Meridian beobachten ließe. An unserer Sternwarte ließ der bedeckte Himmel jedoch vom 19. bis zum 23. März keinen Versuch zu. Am 18. und am 24. dagegen, als der Himmel um die
\SourceLeaf{444}%
Kulminationszeit des Kometen klar war, ließ sich keine Spur des Kometen erkennen, obwohl seine Position genau vorausberechnet worden war. Nach den Nachrichten, die Hofrat \textsc{Gauss} von verschiedenen anderen Sternwarten erhalten hat, sind die Suchversuche dort um die Zeit der größten Helligkeit des Kometen ebenfalls durch ungünstiges Wetter vereitelt worden.
